\section*{Problema 15 del Capítulo 5}

\textbf{Enunciado:}

Calcular los límites siguientes en función del número $\alpha = \lim_{x \to 0} \frac{\sin x}{x}$.

\begin{enumerate}
    \item $\lim_{x \to 0} \frac{\sin 2x}{x}$
    \item $\lim_{x \to 0} \frac{\sin ax}{\sin bx}$
    \item $\lim_{x \to 0} \frac{\sin^2 2x}{x}$
    \item $\lim_{x \to 0} \frac{\sin^2 2x}{x^2}$
    \item $\lim_{x \to 0} \frac{1 - \cos x}{x^2}$
    \item $\lim_{x \to 0} \frac{\tan^2 x + 2x}{x + x^2}$
    \item $\lim_{x \to 0} \frac{x \sin x}{1 - \cos x}$
    \item $\lim_{h \to 0} \frac{\sin(x + h) - \sin x}{h}$
    \item $\lim_{x \to 1} \frac{\sin(x^2 - 1)}{x - 1}$
    \item $\lim_{x \to 0} \frac{x^2(3 + \sin x)}{(x + \sin x)^2}$
    \item $\lim_{x \to 1} (x^2 - 1)^3 \sin\left(\frac{1}{x - 1}\right)^3$
\end{enumerate}

\textbf{Demostración:}

\begin{enumerate}
    \item $\lim_{x \to 0} \frac{\sin 2x}{x} = \lim_{x \to 0} 2 \cdot \frac{\sin 2x}{2x} = 2 \cdot \alpha$.

    \item $\lim_{x \to 0} \frac{\sin ax}{\sin bx} = \lim_{x \to 0} \frac{\sin ax}{ax} \cdot \frac{bx}{\sin bx} \cdot \frac{a}{b} = \frac{a}{b} \cdot \alpha \cdot \frac{1}{\alpha} = \frac{a}{b}$.

    \item $\lim_{x \to 0} \frac{\sin^2 2x}{x} = \lim_{x \to 0} \frac{(\sin 2x)^2}{x} = \lim_{x \to 0} \frac{\sin 2x}{x} \cdot \sin 2x = 2\alpha \cdot 0 = 0$.

    \item $\lim_{x \to 0} \frac{\sin^2 2x}{x^2} = \lim_{x \to 0} \left(\frac{\sin 2x}{x}\right)^2 = (2\alpha)^2 = 4\alpha^2$.

    \item $\lim_{x \to 0} \frac{1 - \cos x}{x^2} = \lim_{x \to 0} \frac{2\sin^2(x/2)}{x^2} = \lim_{x \to 0} \frac{2\left(\frac{\sin(x/2)}{x/2}\right)^2}{4} = \frac{1}{2} \cdot \alpha^2$.

    \item $\lim_{x \to 0} \frac{\tan^2 x + 2x}{x + x^2} = \lim_{x \to 0} \frac{\sin^2 x / \cos^2 x + 2x}{x + x^2} = \lim_{x \to 0} \frac{\sin^2 x}{x} \cdot \frac{1}{x \cos^2 x} + \lim_{x \to 0} \frac{2x}{x + x^2} = \alpha^2 + 2$.

    \item $\lim_{x \to 0} \frac{x \sin x}{1 - \cos x} = \lim_{x \to 0} \frac{x \cdot \frac{\sin x}{x}}{1 - \cos x} = \lim_{x \to 0} \frac{x \alpha}{1 - \cos x} = \lim_{x \to 0} \frac{x \alpha}{2\sin^2(x/2)} = \lim_{x \to 0} \frac{\alpha}{2\sin(x/2)} = \infty$.

    \item $\lim_{h \to 0} \frac{\sin(x + h) - \sin x}{h} = \lim_{h \to 0} \frac{\sin x \cos h + \cos x \sin h - \sin x}{h} = \lim_{h \to 0} \sin x \cdot \frac{\cos h - 1}{h} + \cos x \cdot \frac{\sin h}{h} = \cos x \cdot \alpha$.

    \item $\lim_{x \to 1} \frac{\sin(x^2 - 1)}{x - 1} = \lim_{x \to 1} \frac{\sin(x^2 - 1)}{x^2 - 1} \cdot \frac{x^2 - 1}{x - 1} = \alpha \cdot 2x = 2$.

    \item $\lim_{x \to 0} \frac{x^2(3 + \sin x)}{(x + \sin x)^2} = \lim_{x \to 0} \frac{x^2(3 + \sin x)}{x^2(1 + \sin x / x)^2} = \frac{3 + 0}{(1 + \alpha)^2} = \frac{3}{(1 + \alpha)^2}$.

    \item $\lim_{x \to 1} (x^2 - 1)^3 \sin\left(\frac{1}{x - 1}\right)^3$. Como $\sin\left(\frac{1}{x - 1}\right)$ está acotado entre $-1$ y $1$, se tiene que $(x^2 - 1)^3 \sin\left(\frac{1}{x - 1}\right)^3$ está acotado por $(x^2 - 1)^3$. Cuando $x \to 1$, $(x^2 - 1)^3 \to 0$. Por lo tanto, el límite es $0$.
\end{enumerate}